% !TEX TS-program = pdflatex
% !TeX program = pdflatex
% !TEX encoding = UTF-8
% !TEX spellcheck = fr

\documentclass[xcolor=table]{beamer}


%\usepackage{fullpage}
%\usepackage[left=2.8cm,right=2.2cm,top=2 cm,bottom=2 cm]{geometry}
\setbeamersize{text margin left=10pt,text margin right=10pt}
\usepackage{amsmath,amssymb} 
\usepackage[T1]{fontenc}
\usepackage[utf8]{inputenc}
\usepackage[english,french]{babel}
\usepackage{txfonts}
\usepackage[]{graphicx}
\usepackage{multirow}
\usepackage{hyperref}
\usepackage{colortbl}
\usepackage{listings}
\usepackage{wrapfig}
\usepackage{multicol}
\usepackage[export]{adjustbox} %for images in table, also for frame

\hypersetup{
	colorlinks,
	urlcolor = blue
}

%\renewcommand{\baselinestretch}{1.5}

\def\supit#1{\raisebox{0.8ex}{\small\it #1}\hspace{0.05em}}

\AtBeginSection{%
	\begin{frame}
		\sectionpage
	\end{frame}
}

\newcommand{\rottext}[2]{%
	\rotatebox{90}{%
	\begin{minipage}{#1}%
		\raggedleft#2%
	\end{minipage}%
	}%
}

\usepackage{longtable}
\usepackage{tabu}


\institute{ %
École  nationale Supérieure d'Informatique (ESI, ex. INI), Algérie
}
\author[ \textbf{\footnotesize  \insertframenumber/\inserttotalframenumber} \hspace*{\fill} ESI (2019-2020)] %
{ARIES Abdelkrime}
%\titlegraphic{\includegraphics[height=1cm]{../img/esi-logo.png}%\hspace*{4.75cm}~


\date{Année unniversitaire: 2019/2020} %\today

\usetheme{Warsaw} % Antibes Boadilla Warsaw

\beamertemplatenavigationsymbolsempty

%\setbeamertemplate{headline}{}

\definecolor{lightblue}{HTML}{D0D2FF}
\definecolor{lightyellow}{HTML}{FFFFAA}
\definecolor{darkblue}{HTML}{0000BB}
\definecolor{olivegreen}{HTML}{006600}
\definecolor{violet}{HTML}{6600CC}

\newcommand{\keyword}[1]{\textcolor{red}{\bfseries\itshape #1}}
\newcommand{\expword}[1]{\textcolor{olivegreen}{#1}}
\newcommand{\optword}[1]{\textcolor{violet}{\bfseries #1}}

\makeatletter
\newcommand\mysphere{%
	\parbox[t]{10pt}{\raisebox{0.2pt}{\beamer@usesphere{item projected}{bigsphere}}}}
\makeatother

%\let\oldtabular\tabular
%\let\endoldtabular\endtabular
%\renewenvironment{tabular}{\rowcolors{2}{white}{lightblue}\oldtabular\rowcolor{blue}}{\endoldtabular}


\NoAutoSpacing %french autospacing after ":"

\def\graphpath{}

\newcommand{\changegraphpath}[1]{\def\graphpath{#1}}

\newcommand{\graphintable}[2]{%
	\includegraphics[height=#1,valign=t]{\graphpath #2}%
}

\newcommand{\vgraphpage}[2][.85\textheight]{%
	\begin{center}%
		\includegraphics[height=#1]{\graphpath #2}%
	\end{center}%
}

\newcommand{\hgraphpage}[2][\textwidth]{%
%	\begin{center}%
		\includegraphics[width=#1]{\graphpath #2}%
%	\end{center}%
}

\newcommand{\insertbibliography}[2]{
	\appendix
	\section*{Bibliographie}
	\nocite{#2}
	
	\begin{multicols*}{2}[\frametitle{\insertsection} \usebeamertemplate{frametitle}]%\usebeamertemplate*{frametitle}\frametitle{Références}
		\tiny
		\bibliography{#1}
		\bibliographystyle{acm}
	\end{multicols*}
}


\usepackage{lstautogobble}

\title[BWEB: 03- Rédaction (\LaTeX)] %
{Bureautique et Web \\Chapitre 03: Rédaction d'un document numérique \\ \slshape\small Introduction au \LaTeX}  

\changegraphpath{../img/Bweb03-redaction/latex/}


\begin{document}

\begin{frame}
\frametitle{Rédaction d'un document numérique}
\framesubtitle{\LaTeX}

\begin{itemize}
	\item un formateur de texte ; pas un logiciel de traitement de texte comme Word et Writer
	\item très utilisé dans les articles de recherche (surtout ceux contenant beaucoup d'équations)
	\item il sépare le contenu du format: on peut facilement appliquer des modèles sur un même document
	\item un fichier \LaTeX\ est un fichier texte simple ayant l'extension \optword{.tex}
	\item ce fichier contient des commandes et peut être lu par n'importe quel éditeur de texte
	\item les compilateurs \LaTeX, comme \expword{pdflatex} ou \expword{xelatex}, génèrent des PDF
	
\end{itemize}

\end{frame}

\begin{frame}
\frametitle{Rédaction d'un document numérique}
\framesubtitle{\LaTeX : Plan}

\begin{multicols}{2}
%	\small
\tableofcontents
\end{multicols}
\end{frame}

%===================================================================================
\section{Préparer un document}
%===================================================================================

\begin{frame}[fragile]
\frametitle{Préparer un document}
Un document \LaTeX se compose de ces parties :
\begin{itemize}
	\item préambule, contenant la classe du document (comme un modèle), appels des extensions et des commandes
	\item contenu entre \keyword{\textbackslash begin\{document\}} et \keyword{\textbackslash end\{document\}}
\end{itemize}

\begin{minipage}{0.60\textwidth} 
	Parmi les classes disponibles
	\begin{itemize}
		\item \optword{book} : pour rédiger des livres
%		\item memoir
		\item \optword{report} : pour rédiger des rapports
		\item \optword{article} : pour rédiger des articles
		\item \optword{letter} : pour rédiger des lettres
		\item \optword{beamer} : pour créer des présentations (comme celle-ci)
	\end{itemize}
\end{minipage}
%
\begin{minipage}{0.38\textwidth}
\begin{exampleblock}{exemple: hello.tex}
\scriptsize\bfseries
\begin{lstlisting}[language={[LaTeX]TeX}]
\documentclass{book}
% espace pour les packages
\begin{document}
    % le contenu
    Hello World!
\end{document}
\end{lstlisting}
\end{exampleblock}
\end{minipage}

\end{frame}

\subsection{Sections}

\begin{frame}
\frametitle{Préparer un document}
\framesubtitle{Sections}

\rowcolors{2}{lightblue}{lightyellow}

\begin{tabular}{p{.3\textwidth}cp{.6\textwidth}}
	%	\hline\hline
	\rowcolor{darkblue}
	\textcolor{white}{Section} && \textcolor{white}{Fonction} \\
	%	\hline\hline
	\textbackslash part\{...\} && La partie; ex. \expword{\'Etude bibliographique}. 
	
	Elle n'existe pas dans \optword{letter}.\\
	
	\textbackslash chapter\{...\} && Un chapitre; 
	
	Elle existe seulement dans \optword{book} et \optword{report}.\\
	
	\textbackslash section\{...\} && Une section (un titre de niveau 1); 
	
	Elle n'existe pas dans \optword{letter}.\\
	
	\textbackslash subsection\{...\} && Une sous-section (un titre de niveau 2); 
	
	Elle n'existe pas dans \optword{letter}.\\
	
	\textbackslash subsubsection\{...\} && Une subdivision (un titre de niveau 3); 
	
	Elle n'existe pas dans \optword{letter}.\\
\end{tabular}

Si on ajoute un étoile, la section ne sera pas numérotée. 
Par exemple, \expword{\textbackslash chapter*\{Introduction\}}
\end{frame}

%\begin{frame}[fragile]
%\frametitle{Préparer un document: Sections}
%\framesubtitle{Parties et Chapitres}
%
%\begin{exampleblock}{exemple: titre.tex}
%\scriptsize\bfseries
%\begin{lstlisting}
%\documentclass{book}
%
%\title{Introduction LaTeX}
%\author{Abdelkrime Aries}
%\date{2020}
%
%\begin{document}
%
%    \maketitle
%	\chapter*{Introduction}
%	mot mot mot mot
%
%    \chapter{Un chapite numerote}
%    mot mot mot mot
%    
%    \section{une section}
%	
%    \chapter{Un chapite numerote}
%    mot mot mot mot
%    
%    
%
%\end{document}
%\end{lstlisting}
%\end{exampleblock}
%
%\end{frame}

\subsection{Pages}

%\begin{frame}
%\frametitle{Préparer un document}
%\framesubtitle{Pages}
%
%\end{frame}

\begin{frame}[fragile]
\frametitle{Préparer un document: Pages}
\framesubtitle{Titre (Page de garde)}

\begin{minipage}{0.59\textwidth} 
	\begin{itemize}
		\item On doit introduire le titre, le nom de l'auteur, la date et d'autres informations selon la classe utilisée. 
		\item Pour générer le titre à partir de ces informations, on utilise la commande \keyword{\textbackslash maketitle}
		\item Dans \optword{book} et \optword{report}, \LaTeX\ crée une page de garde (une nouvelle page avec les informations). 
		Dans \optword{article}, \LaTeX\ insère ces informations en première page suivies par un texte 
	\end{itemize}
\end{minipage}
%
\begin{minipage}{0.40\textwidth}
\begin{exampleblock}{exemple: titre.tex}
\scriptsize\bfseries
\begin{lstlisting}
\documentclass{article}
	
\title{Introduction LaTeX}
\author{Abdelkrime Aries}
\date{2020}
	
\begin{document}
	
    \maketitle
	
    Hello World!
	
\end{document}
\end{lstlisting}
\end{exampleblock}
\end{minipage}

\end{frame}

\begin{frame}[fragile]
\frametitle{Préparer un document: Pages}
\framesubtitle{Taille et format}

\begin{itemize}
	\item Dans les options de la classe, on peut définir la taille de la page. 
	Parmi les tailles, on peut citer: \optword{a4paper}, \optword{a5paper} et \optword{letterpaper}.
	\item On peut, aussi, définir la taille des caractères pour le texte normal. 
	Dans les rapports de l'ESI, la taille utilisée est \optword{12pt}.
	\item Si on veut que le document soit recto-verso, on utilise l'option \optword{twoside}
\end{itemize}

\begin{exampleblock}{exemple: la taille de la page}
\scriptsize\bfseries
\begin{lstlisting}
\documentclass[a4paper, 12pt, twoside]{book}
\end{lstlisting}
\end{exampleblock}

\end{frame}

\begin{frame}[fragile]
\frametitle{Préparer un document: Pages}
\framesubtitle{Orientation (tout le document)}

\begin{itemize}
	\item Par défaut, l'orientation \optword{Portrait}
	\item Pour utiliser l'orientation \optword{Paysage}, on fait appel à l'extension \optword{geometry} avec l'option \optword{landscape}
\end{itemize}

\begin{exampleblock}{exemple: orientation paysage}
\scriptsize\bfseries
\begin{lstlisting}
\documentclass[a4paper, 12pt]{book}
\usepackage[landscape]{geometry}
\end{lstlisting}
\end{exampleblock}

\end{frame}


\begin{frame}[fragile]
\frametitle{Préparer un document: Pages}
\framesubtitle{Orientation (une partie du document)}

\begin{minipage}{0.30\textwidth} 
\begin{itemize}
	\item Importer l'extension \optword{pdflscape}
	\item Mettre le contenu dans l'environnement \optword{landscape}
\end{itemize}
\end{minipage}
%
\begin{minipage}{0.69\textwidth}
\begin{exampleblock}{exemple: paysage.tex}
\scriptsize\bfseries
\begin{lstlisting}
\documentclass[a4paper, 12pt]{article}
\usepackage{pdflscape}

\begin{document}
    page 1 avec orientation portrait

    \begin{landscape}
        page 2 avec orientation paysage 
        \newpage
        page 3 avec orientation paysage
    \end{landscape}

    page 4 avec orientation portrait
\end{document}
\end{lstlisting}
\end{exampleblock}
\end{minipage}

\end{frame}


\begin{frame}[fragile]
\frametitle{Préparer un document: Pages}
\framesubtitle{Marges}

\begin{itemize}
	\item On définit les marges en utilisant l'extension \optword{geometry}.
	\item Dans ces options, on peut fixer les marges gauche (\optword{left}), droite (\optword{right}), tête (\optword{top}) et pied (\optword{bottom}).
	\item Si on veut appliquer la même marge pour les quatre coins: \optword{margin}
	\item Si le document est recto-verso, on utilise \optword{inner} pour la marge intérieure (envers le reliure) et \optword{outer} pour la marge extérieure.
\end{itemize}

\begin{exampleblock}{exemple: les marges de la page}
\scriptsize\bfseries
\begin{lstlisting}
\documentclass[a4paper, 12pt]{book}
\usepackage[left=2.8cm, right=2.2cm, top=2cm, bottom=2cm]{geometry}
\end{lstlisting}
\end{exampleblock}

\end{frame}

\begin{frame}[fragile]
\frametitle{Préparer un document: Pages}
\framesubtitle{Colonnes (tout le document)}

\begin{itemize}
	\item Il existe des articles avec deux colonnes
	\item Utiliser l'option \optword{twocolumn} de la classe.
\end{itemize}

\begin{exampleblock}{exemple: les marges de la page}
\scriptsize\bfseries
\begin{lstlisting}
\documentclass[a4paper, 12pt, twocolumn]{article}
\end{lstlisting}
\end{exampleblock}

\end{frame}

\begin{frame}[fragile]
\frametitle{Préparer un document: Pages}
\framesubtitle{Colonnes (une partie du document)}

\begin{minipage}{0.30\textwidth} 
\begin{itemize}
	\item Importer l'extension \optword{multicol}
	\item Mettre le texte dans l'environnement \optword{multicols} en spécifiant le nombre des colonnes
\end{itemize}
\end{minipage}
%
\begin{minipage}{0.69\textwidth}
\begin{exampleblock}{exemple: colonnes.tex}
\scriptsize\bfseries
\begin{lstlisting}
\documentclass{article}
\usepackage{multicol}
	
\begin{document}
\begin{multicols}{3}
[Un texte qui s'affiche sur une seule colonne.]
Ceci est la premiere phrase. Ceci est la phrase 2.
Ceci est la phrase 3. Ceci est la phrase 4.
Ceci est la phrase 5. Ceci est la phrase 6.
Ceci est la phrase 7. Ceci est la phrase 8.
Ceci est la phrase 9. Ceci est la phrase 10.
\end{multicols}
Un texte qui s'affiche sur une seule colonne.
	
\end{document}
\end{lstlisting}
\end{exampleblock}
\end{minipage}

\end{frame}

\begin{frame}[fragile]
\frametitle{Préparer un document: Pages}
\framesubtitle{Saute de page}

Pour insérer un saut de page
\begin{itemize}
	\item on utilise la commande \keyword{\textbackslash newpage}
	\item avant un chapitre, on utilise la commande \keyword{\textbackslash clearpage}
	\item avant un chapitre dans un document recto-verso, on utilise la commande \keyword{\textbackslash cleardoublepage}
\end{itemize}

\begin{exampleblock}{exemple: nouvpage.tex}
\scriptsize\bfseries
\begin{lstlisting}
\documentclass{article}
\begin{document}
    Avant la nouvelle page
    \newpage
    Apres la nouvelle page
\end{document}
\end{lstlisting}
\end{exampleblock}

\end{frame}

%\begin{frame}
%\frametitle{Préparer un document: Pages}
%\framesubtitle{Bordures}
%
%\end{frame}

%\begin{frame}
%\frametitle{Préparer un document: Pages}
%\framesubtitle{Filigrane}
%
%\end{frame}

\begin{frame}[fragile]
\frametitle{Préparer un document: Pages}
\framesubtitle{Arrière-plan}

\begin{minipage}{0.49\textwidth}
Pour définir une couleur d'arrière plan
\begin{itemize}
	\item on utilise la commande \keyword{\textbackslash pagecolor} de l'extension \optword{xcolor}
	\item pour utiliser une image comme arrière-plan, on fait appel à l'extension \optword{background} (voir le document \expword{pageimage.tex})
\end{itemize}
\end{minipage}
%
\begin{minipage}{0.50\textwidth}
\begin{exampleblock}{exemple: pagecouleur.tex}
\scriptsize\bfseries
\begin{lstlisting}
\documentclass{article}
\usepackage{xcolor}
	
\begin{document}
    \pagecolor{yellow}
    Je suis une page jaune

    \newpage

    \pagecolor{green}
    Je suis une page verte
	
\end{document}
\end{lstlisting}
\end{exampleblock}
\end{minipage}

\end{frame}

\begin{frame}
\frametitle{Préparer un document: Pages}
\framesubtitle{En-tête et pied de page: fancyhdr}


\begin{itemize}
	\item utiliser l'extension \optword{fancyhdr}
	\item utiliser la commande \keyword{\textbackslash pagestyle}:
	\begin{itemize}
		\item \optword{empty} l'entête et le pied sont vides
		\item \optword{fancy} entête et le pied de \optword{fancyhdr}
	\end{itemize}
	\item configurer l'entête par \keyword{\textbackslash fancyhead} et le pied par \keyword{\textbackslash fancyfoot}
	\item la syntaxe est \expword{\textbackslash fancyhead[sélecteurs]\{le texte à écrire\}}
	\item les sélecteurs sont
	\begin{itemize}
		\item \optword{E} page paire, \optword{O} page impaire 
		\item \optword{L} gauche, \optword{C} centre, \optword{R} droit 
	\end{itemize}
	\item utiliser \expword{\textbackslash fancyhf\{\}} pour remettre à zéro
	\item exemple \expword{\textbackslash fancyhead[LE]\{Intro à LaTeX\}}: écrire ce titre dans l'entête gauche des pages paires
\end{itemize}

\end{frame}

\begin{frame}[fragile]
\frametitle{Préparer un document: Pages}
\framesubtitle{En-tête et pied de page: exemple}

\begin{exampleblock}{exemple: paragraphes.tex}
\scriptsize\bfseries\vspace{-6pt}
\begin{lstlisting}
\documentclass{book}
\usepackage{fancyhdr}

\renewcommand{\headrulewidth}{0.5pt}
\renewcommand{\footrulewidth}{2pt}

\fancyhf{}
\fancyhead[R]{ARIES Abdelkrime}
\fancyhead[LE]{\leftmark}
\fancyhead[LO]{\rightmark}
\fancyfoot[RO, LE]{\thepage}

\begin{document}
    ...
    \pagestyle{fancy}
    ...
\end{document}
\end{lstlisting}\vspace{-6pt}
\end{exampleblock}


\end{frame}

\subsection{Paragraphes}

\begin{frame}[fragile]
\frametitle{Préparer un document}
\framesubtitle{Paragraphes}

\begin{minipage}{0.49\textwidth}
	\begin{itemize}
		\item Si on laisse une ligne vide, \LaTeX créera un nouveau paragraphe. 
		\item On peut écrire les phrases d'un même paragraphe sur multiples lignes.
		\item Pour sauter la ligne sans créer un nouveau paragraphe, utiliser la commande \keyword{\textbackslash newline} ou simplement \keyword{\textbackslash \textbackslash}
	\end{itemize}
\end{minipage}
%
\begin{minipage}{0.50\textwidth}
\begin{exampleblock}{exemple: paragraphes.tex}
\scriptsize\bfseries
\begin{lstlisting}
\documentclass{article}
	
\begin{document}
    Phrase 1 Paragraphe 1. 
    Phrase 2 Paragraphe 1. 
    
    Phrase 1 Paragraphe 2. 
    Phrase 2 Paragraphe 2.	
    \newline	
    Phrase 3 Paragraphe 2.
\end{document}
\end{lstlisting}
\end{exampleblock}
\end{minipage}

\end{frame}


\begin{frame}
\frametitle{Préparer un document: Paragraphes}
\framesubtitle{Internationalisation}

\begin{itemize}
	\item pour rédiger avec d'autres langues, on utilise l'extension \optword{babel} avec une liste des langues comme option
	\item la dernière langue dans les options est celle activée au départ
	\item pour utiliser l'encodage \keyword{UTF-8} comme méthode d'entrée, on utilise l'extension \optword{inputenc}
	\item pour sélectionner une langue, on utilise la commande \keyword{\textbackslash selectlanguage}
	\item pour écrire un texte avec une autre langue sans changer celle activée, on utilise l'environnement \optword{otherlanguage}
\end{itemize}

\end{frame}

\begin{frame}[fragile]
\frametitle{Préparer un document: Paragraphes}
\framesubtitle{Internationalisation (exemple)}

\begin{exampleblock}{exemple: i18n.tex}
\lstset{escapeinside=**}\novocalize
\scriptsize\bfseries\vspace{-6pt}
\begin{lstlisting}
\documentclass{article}
\usepackage[utf8]{inputenc}
\usepackage[arabic, french]{babel}

\begin{document}
    *Les caractères é, è, à, ô, etc. ne causeront aucun problème. *
    \selectlanguage{arabic}
    *\<h_dA n.s bAl`rbyT>*
    \selectlanguage{french}
    *Un texte en français avec un autre en arabe*

    \begin{otherlanguage}{arabic}
       *\<h_dA n.s bAl`rbyT>*
    \end{otherlanguage}
\end{document}
\end{lstlisting}\vspace{-6pt}
\end{exampleblock}

\end{frame}

\begin{frame}[fragile]
\frametitle{Préparer un document: Paragraphes}
\framesubtitle{Alignement}

\begin{minipage}{0.49\textwidth}
	Les environnements:
	\begin{itemize}
		\item \optword{flushleft} pour aligner à gauche
		\item \optword{center} pour centrer
		\item \optword{flushright} pour aligner à droite
	\end{itemize}
\end{minipage}
%
\begin{minipage}{0.50\textwidth}
\begin{exampleblock}{exemple: alignement.tex}
\scriptsize\bfseries
\begin{lstlisting}
\documentclass{article}

\begin{document}
    \begin{flushleft}
        Aligner a gauche
    \end{flushleft}

    \begin{center}
        Centrer
    \end{center} 

    \begin{flushright}
        Aligner a droite
    \end{flushright}
\end{document}
\end{lstlisting}
\end{exampleblock}
\end{minipage}

\end{frame}


\begin{frame}[fragile]
\frametitle{Préparer un document: Paragraphes}
\framesubtitle{Espacement vertical, interligne et retrait de 1ère ligne}

\begin{minipage}{0.49\textwidth}
	Les paramètres :
	\begin{itemize}
		\item \keyword{\textbackslash parskip} : espacement entre deux paragraphes
		\item \keyword{\textbackslash parindent} : retrait de 1ère ligne
		\item \keyword{\textbackslash baselineskip} : interligne
	\end{itemize}

	Avant un paragraphe :
	\begin{itemize}
		\item \keyword{\textbackslash indent} : forcer le retrait
		\item \keyword{\textbackslash noindent} : ignorer le retrait
	\end{itemize}
\end{minipage}
%
\begin{minipage}{0.50\textwidth}
\begin{exampleblock}{exemple: espacement.tex}
\scriptsize\bfseries
\begin{lstlisting}
\documentclass{article}

\setlength{\parskip}{1cm}
\setlength{\baselineskip}{12pt}
\setlength{\parindent}{2cm}

\begin{document}

    Phrase 1 Paragraphe 1. \newline
    Phrase 2 Paragraphe 1. 

    Phrase 1 Paragraphe 2. 
    Phrase 2 Paragraphe 2.		

\end{document}
\end{lstlisting}
\end{exampleblock}
\end{minipage}

\end{frame}

\begin{frame}[fragile]
\frametitle{Préparer un document: Paragraphes}
\framesubtitle{Listes}

\begin{minipage}{0.49\textwidth}
	Les environnements :
	\begin{itemize}
		\item \optword{itemize} : listes à puces
		\item \optword{enumerate} : listes numérotées
	\end{itemize}
	
	caractérisés par :
	\begin{itemize}
		\item un élément est défini par \keyword{\textbackslash item}
		\item on peut créer des sous-listes
	\end{itemize}
\end{minipage}
%
\begin{minipage}{0.50\textwidth}
\begin{exampleblock}{exemple: listes.tex}
\scriptsize\bfseries
\begin{lstlisting}
\documentclass{article}

\begin{document}
    \begin{enumerate}
        \item etape 1
        \item etape 2:
        \begin{itemize}
            \item un element
            \item un autre element
        \end{itemize}
    \end{enumerate}
\end{document}
\end{lstlisting}
\end{exampleblock}
\end{minipage}

\end{frame}


\subsection{Mise en forme}

\begin{frame}[fragile]
\frametitle{Préparer un document: Mise en forme}
\framesubtitle{Styles de la police}

\begin{minipage}{0.49\textwidth}
	Les commandes :
	\begin{itemize}
		\item \keyword{\textbackslash textbf} : gras
		\item \keyword{\textbackslash textit} : italique
		\item \keyword{\textbackslash underline} : souligné
		\item \keyword{\textbackslash uppercase} : majuscule
		\item \keyword{\textbackslash lowercase} : minuscule
		\item \keyword{\textbackslash textrm} : famille roman 
		\item \keyword{\textbackslash textsf} : famille sans serif
		\item \keyword{\textbackslash texttt} : famille teletype
	\end{itemize}
\end{minipage}
%
\begin{minipage}{0.50\textwidth}
\begin{exampleblock}{exemple: styles.tex}
\scriptsize\bfseries
\begin{lstlisting}
\documentclass{article}
		
\begin{document}
    \textbf{Texte En Gras}. 
    \textit{Texte En Italique}. 
    \underline{Texte Souligne}. 
    \uppercase{Texte En Majuscule}. 
    \lowercase{Texte En Minuscule}.
    \textrm{texte en roman}. 
    \textsf{texte en sans serif}. 
    \texttt{texte en teletype}. 
\end{document}
\end{lstlisting}
\end{exampleblock}
\end{minipage}


\end{frame}

\begin{frame}[fragile]
\frametitle{Préparer un document: Mise en forme}
\framesubtitle{Tailles de la police}

\begin{minipage}{0.39\textwidth}
	\begin{itemize}
		\item \keyword{\textbackslash tiny}
		\item \keyword{\textbackslash scriptsize}
		\item \keyword{\textbackslash footnotesize}
		\item \keyword{\textbackslash small}
		\item \keyword{\textbackslash normalsize}
		\item \keyword{\textbackslash large}
		\item \keyword{\textbackslash Large}
		\item \keyword{\textbackslash LARGE}
		\item \keyword{\textbackslash huge}
		\item \keyword{\textbackslash Huge}
	\end{itemize}
\end{minipage}
%
\begin{minipage}{0.60\textwidth}
\begin{exampleblock}{exemple: taillespolice.tex}
\scriptsize\bfseries
\begin{lstlisting}
\documentclass[12pt]{article}
		
\begin{document}
    {\tiny texte en tiny.} 
    {\scriptsize texte en scriptsize.} 
    {\footnotesize texte en footnotesize.} 
    {\small texte en small.} 
    {\normalsize texte en normalsize.} 
    {\large texte en large.} 
    {\Large texte en Large.} 
    {\LARGE texte en LARGE.} 
    {\huge texte en huge.} 
    {\Huge texte en Huge.} 
\end{document}
\end{lstlisting}
\end{exampleblock}
\end{minipage}

\end{frame}

\begin{frame}[fragile]
\frametitle{Préparer un document: Mise en forme}
\framesubtitle{Couleur du texte}

\begin{minipage}{0.39\textwidth}
	\begin{itemize}
		\item utiliser la commande \keyword{\textbackslash textcolor} de l'extension \optword{xcolor}
		\item les couleurs prédéfinies sont: black, blue, brown, cyan, darkgray, gray, green, lightgray, lime, magenta, olive, orange, pink, purple, red, teal, violet, white, yellow.
		\item pour définir d'autres couleurs, utiliser la commande \keyword{\textbackslash definecolor}
	\end{itemize}
\end{minipage}
%
\begin{minipage}{0.60\textwidth}
\begin{exampleblock}{exemple: taillespolice.tex}
\scriptsize\bfseries
\begin{lstlisting}
\documentclass{article}
\usepackage{xcolor}
\definecolor{couleur1}{RGB}{180, 80, 255}
\begin{document}

    Un texte normal. 

    \textcolor{red}{Un texte en rouge}

    \textcolor{couleur1}{Un texte en couleur1}

\end{document}
\end{lstlisting}
\end{exampleblock}
\end{minipage}

\end{frame}


%===================================================================================
\section{Enrichir un document}
%===================================================================================

%\begin{frame}
%\frametitle{Enrichir un document}
%
%\end{frame}

\subsection{Tableaux}

\begin{frame}[fragile]
\frametitle{Enrichir un document : Tableaux}
\framesubtitle{Insérer un tableau}

\begin{minipage}{0.59\textwidth}
	\begin{itemize}
		\item utiliser l'environnement \optword{tabular}
		\item l'argument définit le nombre des colonnes en utilisant les lettres: \optword{l} (aligner à gauche), \optword{c} (center) et \optword{r} (aligner à droite).
		\item la barre (\optword{|}) trace un trait vertical  
		\item \keyword{\textbackslash hline} trace un trait horizontal
		\item les colonnes sont séparés par \optword{\&} et les lignes par \optword{\textbackslash \textbackslash}
	\end{itemize}
\end{minipage}
%
\begin{minipage}{0.40\textwidth}
\begin{exampleblock}{exemple: tableaux1.tex}
\scriptsize\bfseries
\begin{lstlisting}
\documentclass{article}

\begin{document}
    \begin{tabular}{|l|c|r|}
        \hline
        l1c1 & l1c2 & l1c3 \\
        \hline 
        l2c1 & l2c2 & l2c3 \\
        \hline
    \end{tabular}
\end{document}
\end{lstlisting}
\end{exampleblock}
\end{minipage}

\end{frame}


\begin{frame}[fragile]
\frametitle{Enrichir un document : Tableaux}
\framesubtitle{Positionner un tableau}

\begin{minipage}{0.59\textwidth}
	\begin{itemize}
		\item utiliser l'environnement \optword{table}
		\item l'argument définit le positionnement en utilisant les lettres:
		\begin{itemize}
			\item \optword{!} demande à \LaTeX\ de faire tout son possible pour respecter l'ordre indiqué
			ensuite ;
			\item \optword{h} (here) ici, si possible ;
			\item \optword{t} (top) en haut d’une page (celle-ci ou la suivante) ;
			\item \optword{b} (bottom) en bas d’une page (celle-ci ou la suivante) ;
			\item \optword{p} (page of floats) sur une page spéciale ne contenant pas de texte mais
			uniquement des tableaux et des figures.
			
		\end{itemize}
		
	\end{itemize}
\end{minipage}
%
\begin{minipage}{0.40\textwidth}
\begin{exampleblock}{exemple: tableaux2.tex}
\scriptsize\bfseries
\begin{lstlisting}
\documentclass{article}
		
\begin{document}
   \begin{table}[!htp]
      \begin{tabular}{|l|c|r|}
         \hline
         l1c1 & l1c2 & l1c3 \\
         \hline 
         l2c1 & l2c2 & l2c3 \\
         \hline
      \end{tabular}
   \end{table}
\end{document}
\end{lstlisting}
\end{exampleblock}
\end{minipage}

\end{frame}

\subsection{Images}

\begin{frame}[fragile]
\frametitle{Enrichir un document: Images}
\framesubtitle{Insérer une image}

\begin{itemize}
	\item utiliser la commande \keyword{\textbackslash includegraphics} de l'extension \optword{graphicx}
	\item on peut fixer la largeur (\optword{width}) et la hauteur (\optword{height})
	\item pour récupérer la largeur du texte, utiliser la commande \keyword{\textbackslash textwidth}
	\item pour récupérer la hauteur du texte, utiliser la commande \keyword{\textbackslash textheight}
\end{itemize}

\begin{exampleblock}{exemple: image1.tex}
\scriptsize\bfseries
\begin{lstlisting}
\documentclass{article}
\usepackage{graphicx}
\begin{document}
    \includegraphics[width=4cm, height=2cm]{banner.jpg}
\end{document}
\end{lstlisting}
\end{exampleblock}

\end{frame}

\begin{frame}[fragile]
\frametitle{Enrichir un document: Images}
\framesubtitle{Positionner une figure}

\begin{minipage}{0.54\textwidth}
	\begin{itemize}
		\item utiliser l'environnement \optword{figure}
		\item l'argument définit le positionnement en utilisant les lettres:
		\begin{itemize}
			\item \optword{!} demande à \LaTeX\ de faire tout son possible pour respecter l'ordre indiqué
			ensuite ;
			\item \optword{h} (here) ici, si possible ;
			\item \optword{t} (top) en haut d’une page (celle-ci ou la suivante) ;
			\item \optword{b} (bottom) en bas d’une page (celle-ci ou la suivante) ;
			\item \optword{p} (page of floats) sur une page spéciale ne contenant pas de texte mais
			uniquement des tableaux et des figures.
			
		\end{itemize}
		
	\end{itemize}
\end{minipage}
%
\begin{minipage}{0.45\textwidth}
\begin{exampleblock}{exemple: image2.tex}
\scriptsize\bfseries
\begin{lstlisting}
\documentclass{article}
\usepackage{graphicx}

\begin{document}
   \begin{figure}[!htp]
      \includegraphics{banner.jpg}
   \end{figure}
\end{document}
\end{lstlisting}
\end{exampleblock}
\end{minipage}

\end{frame}

\subsection{Liens et renvoi}

\begin{frame}[fragile]
\frametitle{Enrichir un document: Liens et renvoi}
\framesubtitle{Insérer un lien hypertexte}

\begin{itemize}
	\item importer l'extension \optword{hyperref}
	\item utiliser la commande \keyword{\textbackslash url} pour afficher un URL
	\item utiliser la commande \keyword{\textbackslash href} pour afficher une description avec un lien vers un URL
\end{itemize}

\begin{exampleblock}{exemple: liens.tex}
\scriptsize\bfseries
\begin{lstlisting}
\documentclass{article}
\usepackage{hyperref}
		
\begin{document}
   Le site de l'ESI : \url{http://esi.dz/}
   
   \href{http://esi.dz/}{Cliquer ici pour ouvrir le site de l'ESI}
\end{document}
\end{lstlisting}
\end{exampleblock}


\end{frame}

\begin{frame}[fragile]
\frametitle{Enrichir un document: Liens et renvoi}
\framesubtitle{Renvoi vers un emplacement dans le document}

\begin{itemize}
	\item importer l'extension \optword{hyperref}
	\item utiliser la commande \keyword{\textbackslash hypertarget} pour marquer la destination
	\item utiliser la commande \keyword{\textbackslash hyperlink} pour ajouter un lien vers la destination
\end{itemize}

\begin{exampleblock}{exemple: liens.tex}
\scriptsize\bfseries
\begin{lstlisting}
\documentclass{article}
\usepackage{hyperref}

\begin{document}
    Aller vers \hyperlink{mapage2}{la page suivante}
    \newpage
    \hypertarget{mapage2}{Ici} c'est la destination
\end{document}
\end{lstlisting}
\end{exampleblock}

\end{frame}

\subsection{Formules et symboles}

\begin{frame}[fragile]
\frametitle{Enrichir un document: Formules et symboles}
\framesubtitle{Insérer une formule}

\begin{minipage}{0.45\textwidth}
	Trois types de formules: 
	\begin{itemize}
		\item avec le texte: \optword{\$<formule>\$}
		\item nouvelle ligne sans numérotation: \optword{\textbackslash [<formule>\textbackslash ]}
		\item environnements: le plus utilisé est \optword{equation}
		\item pour capter le numéro d'une équation, on l'affecte une étiquette en utilisant la commande \keyword{\textbackslash label} 
		\item pour faire référence à cette équation, on utilise la commande \keyword{\textbackslash ref} 
	\end{itemize}
\end{minipage}
%
\begin{minipage}{0.54\textwidth}
\begin{exampleblock}{exemple: equation1.tex}
\scriptsize\bfseries
\begin{lstlisting}
\documentclass{article}
\begin{document}
   La probabilite conditionnelle 
   $ P(A/B) = \frac{P(A \cap B)}{P(B)} $
   ce qui veut dire 
   \[ P(A \cap B) = P(A/B) P(B) \]
   Si A et B sont independants, 
   cette probabilite se calcule 
   par l'equation \ref{eq-indep}
   \begin{equation}
      \label{eq-indep}
      P(A \cap B) = P(A) P(B)
   \end{equation}
\end{document}
\end{lstlisting}
\end{exampleblock}
\end{minipage}

\end{frame}

\begin{frame}
\frametitle{Enrichir un document: Formules et symboles}
\framesubtitle{Symboles mathématiques}

symboles d'opérateurs

\keyword{\textbackslash times} ($ \times $), 
\keyword{\textbackslash div} ($ \div $), 
\keyword{\textbackslash cup} ($ \cup $), 
\keyword{\textbackslash cap} ($ \cap $), 
\keyword{\textbackslash wedge} ($ \wedge $), 
\keyword{\textbackslash vee} ($ \vee $ ),
\keyword{\textbackslash prod} ($ \prod $),
\keyword{\textbackslash sum} ($ \sum $),
\keyword{\textbackslash int} ($ \int $)

\vfill

symboles de relation

\keyword{\textbackslash propto} ($ \propto $), 
\keyword{\textbackslash equiv} ($ \equiv $), 
\keyword{\textbackslash neq} ($ \neq $), 
\keyword{\textbackslash approx} ($ \approx $), 
\keyword{\textbackslash sim} ($ \sim $), 
\keyword{\textbackslash perp} ($ \perp $ ),
\keyword{\textbackslash geq} ($ \geq $),
\keyword{\textbackslash leq } ($ \leq  $),
\keyword{\textbackslash subset } ($ \subset$),
\keyword{\textbackslash subseteq } ($\subseteq$),
\keyword{\textbackslash supset } ($ \supset$),
\keyword{\textbackslash supseteq } ($ \supseteq$)

\vfill

symboles de flèches

\keyword{\textbackslash rightarrow} ($ \rightarrow $), 
\keyword{\textbackslash leftarrow} ($ \leftarrow $),
\keyword{\textbackslash leftrightarrow} ($ \leftrightarrow $),
\keyword{\textbackslash Longrightarrow} ($ \Longrightarrow $),
\keyword{\textbackslash Longleftarrow} ($ \Longrightarrow $),
\keyword{\textbackslash Longleftrightarrow} ($ \Longleftrightarrow $)

\vfill

Lettres grecques 

\keyword{\textbackslash alpha} ($ \alpha $), 
\keyword{\textbackslash beta} ($ \beta $),
\keyword{\textbackslash gamma} ($ \gamma $),
\keyword{\textbackslash delta} ($ \delta $),
\keyword{\textbackslash epsilon} ($ \epsilon $),
\keyword{\textbackslash varepsilon} ($ \varepsilon $),
\keyword{\textbackslash theta} ($ \theta $),
\keyword{\textbackslash kappa} ($ \kappa $),
\keyword{\textbackslash pi} ($ \pi $),
\keyword{\textbackslash sigma} ($ \sigma $)
 
\end{frame}


\begin{frame}[fragile]
\frametitle{Enrichir un document: Formules et symboles}
\framesubtitle{Quelques équations mathématiques}

Pour avoir d'autres symboles et fonctions, importer les extensions \optword{amsmath} et \optword{amssymb}

\vfill

{\scriptsize \verb|\int_{0}^{\pi}\cos(mt)\cos(nt)\mathrm{d}t=0|} \hfill
$ \int_{0}^{\pi}\cos(mt)\cos(nt)\mathrm{d}t=0 $

{\scriptsize \verb|\lim_{z\to 0}\frac{\sin z}{z}=1|} \hfill
$ \lim_{z\to 0}\frac{\sin z}{z}=1 $

{\scriptsize \verb|\lim\limits_{z\to 0}\frac{\sin z}{z}=1|} \hfill
$ \lim\limits_{z\to 0}\frac{\sin z}{z}=1 $

{\scriptsize \verb|i = \sqrt{-1}|} \hfill
$ i = \sqrt{-1} $

{\scriptsize \verb|\log_aN = x \Leftrightarrow N = a^x|} \hfill
$ \log_aN = x \Leftrightarrow N = a^x $
	
\end{frame}

%===================================================================================
\section{Références}
%===================================================================================

%\begin{frame}
%\frametitle{Références}
%
%\end{frame}

\subsection{Notes de bas de page}

\begin{frame}[fragile]
\frametitle{Références : Notes de bas de page}
\framesubtitle{Insérer une note de bas de page}

\begin{itemize}
	\item utiliser la commande \keyword{\textbackslash footnote}
\end{itemize}

\begin{exampleblock}{exemple: footnote.tex}
\scriptsize\bfseries
\begin{lstlisting}
\documentclass{article}
\begin{document}
	Un livre contien une page de garde
	\footnote{page de garde: premiere page dans le livre}. 
	Il peut contenir un index
	\footnote{l'index contient les mots cles}. 
	\newpage
	Il peut contenir un index
	\footnote{l'index contient les mots cles}. 
\end{document}
\end{lstlisting}
\end{exampleblock}

\end{frame}

\subsection{Tables (de matières et des illustrations)}

\begin{frame}[fragile]
\frametitle{Références : Tables (de matières et des illustrations)}
\framesubtitle{Table de matières}

\begin{itemize}
	\item utiliser la commande \keyword{\textbackslash tableofcontents}
\end{itemize}

\begin{exampleblock}{exemple: table-matiere.tex}
\scriptsize\bfseries
\begin{lstlisting}
\documentclass{book}
\begin{document}
    \tableofcontents
    
    \chapter{Microsoft Word} 
    \section*{Introduction} \section{Installer Word} 
    \section{Utiliser Word} \section*{Conclusion}
	
    \chapter{LaTeX} 
    \section*{Introduction} \section{Installer LaTeX} 
    \section{Utiliser LaTeX} \section*{Conclusion}
\end{document}
\end{lstlisting}
\end{exampleblock}

\end{frame}

\begin{frame}[fragile]
\frametitle{Références : Tables (de matières et des illustrations)}
\framesubtitle{Table des illustrations}

\begin{itemize}
	\item dans les environnements \optword{table} et \optword{figure}
	\item ajouter la légende en utilisant la commande \keyword{\textbackslash caption}
	\item marquer l'illustration par une étiquette en utilisant la commande  \keyword{\textbackslash label}
	\item on peut référencer une étiquette en utilisant la commande  \keyword{\textbackslash ref}
	\item pour ajouter la liste des tableaux, utiliser la commande \keyword{\textbackslash listoftables}
	\item pour ajouter la liste des figures, utiliser la commande \keyword{\textbackslash listoffigures}
\end{itemize}

\end{frame}

\begin{frame}[fragile]
\frametitle{Références : Tables (de matières et des illustrations)}
\framesubtitle{Table des illustrations (exemple)}

\vspace{-6pt}
\begin{exampleblock}{exemple: illustrations.tex}
\scriptsize\bfseries
\begin{lstlisting}[multicols=2]
\documentclass{book}
\usepackage{graphicx}
\begin{document}
  \listoffigures
  \listoftables
  \clearpage
  Dans Tableau~\ref{tab-comp}
  \begin{table}[!ht]
    \caption{Comparaison entre A et B}
    \label{tab-comp}
    \centering
    \begin{tabular}{lll}
    ...
    \end{tabular}
  \end{table}

  \clearpage
  Dans Figure~\ref{fig-exemp}, 
  \begin{figure}[!ht]
    \caption{Un exemple d'image}
    \label{fig-exemp}
    \includegraphics{banner.jpg}
  \end{figure}

  \clearpage
  Dans Tableau~\ref{tab-comp2}
  \begin{table}[!ht]
    \caption{Comparaison2}
    \label{tab-comp2}
    ...
  \end{table}

\end{document}
\end{lstlisting}
\end{exampleblock}

\end{frame}


\subsection{Index}

\begin{frame}[fragile]
\frametitle{Références}
\framesubtitle{Insérer un index}

\begin{minipage}{0.45\textwidth}
	\begin{itemize}
		\item importer l'extension \optword{makeidx}
		\item appeler la commande \keyword{\textbackslash makeindex} dans la préambule
		\item marquer une entrée avec \keyword{\textbackslash index}
		\item afficher l'index avec \keyword{\textbackslash printindex} 
		\item si l'index n'apparait pas, essayer d'exécuter \expword{makeindex <nomfichier>} et compiler une deuxième fois
	\end{itemize}
\end{minipage}
%
\begin{minipage}{0.54\textwidth}
\begin{exampleblock}{exemple: index.tex}
\scriptsize\bfseries
\begin{lstlisting}
\documentclass{article}
\usepackage{makeidx}
\makeindex
\begin{document}

  Instruction Writeln\index{writeln} 
  est ... 
  \newpage
  Un programme commence par 
  Program\index{program}
  \clearpage

  \printindex
\end{document}
\end{lstlisting}
\end{exampleblock}
\end{minipage}

\end{frame}


\subsection{Bibliographie}

\begin{frame}[fragile]
\frametitle{Références}
\framesubtitle{Bibliographie}

Il existe deux façons pour introduire une bibliographie :
\begin{itemize}
	\item en utilisant l'environnement \optword{thebibliography}
	\item en utilisant BibTeX
\end{itemize}

On s'intéresse au deuxième, puisque: 
\begin{itemize}
	\item on peut facilement insérer une nouvelle source 
	\item les citations Bibtex peuvent être trouvées sur le web pour un article ou un livre donné
	\item on peut générer d'autres styles de bibliographie
\end{itemize}

\end{frame}

\begin{frame}[fragile]
\frametitle{Références: Bibliographie}
\framesubtitle{Gestion des sources}

\begin{minipage}{0.45\textwidth}
	\begin{itemize}
		\item les sources sont des fichiers textes ayants l'extension \optword{.bib}
		\item quelques types de sources: \optword{book}, \optword{inproceedings}, \optword{article}, \optword{patent}
		\item fournir une étiquette pour référencer la source 
		\item fournir des informations comme le titre, l'année, les auteurs, le lien, etc.
	\end{itemize}

\end{minipage}
%
\begin{minipage}{0.54\textwidth}
\begin{block}{format d'une source}
\scriptsize\bfseries\vspace{-6pt}
\begin{lstlisting}
@type{etiquette,
  author = {Pren1 Nom1 and Pren2 Nom2},
  title = {le titre de l'ouvrage},
  year = {l'anne de publication},
  ... d'autres informations
}
\end{lstlisting}\vspace{-6pt}
\end{block}


\begin{exampleblock}{exemple: citation.bib}
\scriptsize\bfseries\vspace{-6pt}
\begin{lstlisting}
@book{2012-boudoin,
  author = {Marc Baudoin},
  title = {Apprends LaTeX},
  year = {2012}
}
...
\end{lstlisting}\vspace{-6pt}
\end{exampleblock}
\end{minipage}

\end{frame}

\begin{frame}[fragile]
\frametitle{Références: Bibliographie}
\framesubtitle{Citer des sources et insérer la bibliographie}

\begin{minipage}{0.45\textwidth}
	\begin{itemize}
		\item choisir un style en utilisant \keyword{\textbackslash bibliographystyle}
		\item quelques styles: \optword{plain}, \optword{unsrt}, \optword{alpha} et \optword{abbrv} 
		\item pour citer une source, utiliser \keyword{\textbackslash cite} avec son étiquette
		\item pour insérer la bibliographie, utiliser \keyword{\textbackslash bibliography}
	\end{itemize}
\end{minipage}
%
\begin{minipage}{0.54\textwidth}
\begin{exampleblock}{exemple: biblio.tex}
\scriptsize\bfseries
\begin{lstlisting}
\documentclass{article}
\bibliographystyle{abbrv}
\begin{document}

  LaTeX est un formateur 
  de texte \cite{2012-boudoin}.
  Exemple d'un article 
  de recherche en LaTeX : 
  \cite{2015-aries-al}
  
  \bibliography{citation}
\end{document}
\end{lstlisting}
\end{exampleblock}
\end{minipage}

\end{frame}

\begin{frame}
\frametitle{Rédaction d'un document numérique: \LaTeX}
\framesubtitle{Un peu d'humeur}

\begin{center}
	\vgraphpage{latex-humour.png}
\end{center}

\end{frame}

\insertbibliography{Bweb03}{*}


\end{document}

